\section{本质的差异：七组对偶}\label{sec:duality}

上一章论证了两策略在\emph{信号层}的同构。但同构不等于同质。本章给出七组对偶，逐条说明在共享同一信号的前提下，两个方向为什么会在假设、分布、机制与约束上分道扬镳。每组对偶按“命题—论证—证据—工程映射—边界”的顺序展开，并在末尾给出可检验的判断。

\subsection{对偶一：平稳性与非平稳性}

\textbf{命题}：均值回归策略要求被交易的对象具备平稳性（或至少存在协整关系）；趋势策略反而依赖非平稳性。两者对“世界是什么样”的假设正好相反。

\textbf{论证}。均值回归的可交易基础是价差 $X_t$ 的平稳性：只有当 $X_t$ 存在一个吸引子（长期均值）时，“偏离终将修复”才是可检验的命题。若 $X_t$ 含单位根，则偏离不会回归，反转策略面临无限累积的亏损。因此配对交易的现代形式把协整检验当作入场券：Engle \& Granger(1987) 的两步法——对两个价格序列做回归、再对残差做单位根检验——是这一门槛的标准实现\cite{englegranger1987}；Dickey--Fuller 检验则是单位根判断的入口\cite{dickeyfuller1979}。

趋势策略的处境完全不同。若价格服从带正漂移的随机游走，趋势信号仍然有效——因为漂移本身提供了方向。Hurst 等(2017) 的百年样本显示趋势跟踪在极端宏观环境下表现最好，本质上是因为那些时期价格序列呈强非平稳（危机后的持续单向运动）\cite{hop2017}。换言之，\textbf{趋势策略恰恰在反转策略最不该交易的时刻获利}。

\textbf{工程映射}。这一差异直接决定了策略研发的第一步动作不同。反转策略的研发流程必须包含\emph{可行性筛选}：先做 ADF/协整检验，通过者才进入交易规则设计；筛选失败的样本被丢弃。趋势策略没有对应的筛选步骤，它的第一步是估计动量统计量（$t$ 值、胜率、夏普），信号弱但仍在。

\textbf{边界}。这个对偶有一个重要的反例：\emph{短期反转}（日频、周频）并不依赖严格的协整关系，它更接近微观结构效应——买卖价差反弹、流动性提供者的补偿。这类反转的“平稳性”来自订单流机制而非基本面的长期均衡，其半衰期以小时或天计。因此，把“反转必先协整”绝对化会漏掉整个高频反转族。准确的表述是：\emph{中长尺度反转依赖协整，短尺度反转依赖微观结构}。

\textbf{可检验判断}：在同一资产池上，反转策略的绩效对“是否通过平稳性筛选”应当高度敏感；趋势策略的绩效则不应因剔除不平稳样本而系统性改善。

\subsection{对偶二：收益偏度——凸性与凹性}

\textbf{命题}：趋势策略的收益右偏、本质上是持有看涨期权式的凸性暴露；反转策略的收益左偏、本质上是卖出期权式的凹性暴露。二者在小概率极端事件上表现相反。

\textbf{论证}。Fung \& Hsieh(2001) 用“回溯跨式”（lookback straddle）期权组合复制趋势跟踪者的收益，指出其收益结构具有明确的凸性：在市场发生大幅单边运动时收益放大，在震荡市中缓慢失血\cite{funghsieh2001}。这一刻画解释了业界常说的“crisis alpha”：趋势策略在股灾、危机中往往正收益，因为危机会制造持续的单向运动。

反转策略的偏度方向相反。它的收益来自价格向均值收敛，在绝大多数时间里稳定收租，但在价格\emph{不}收敛时遭受集中损失。这在数学上等价于卖出价外期权：收益分布是“多次小赢、偶尔巨亏”。Harvey \& Siddique(2000) 的条件偏度定价文献为这类暴露提供了资产定价层面的框架\cite{harveysiddique2000}。

\textbf{证据}。趋势一侧有两个直接的实证：Daniel \& Moskowitz(2016) 记录的动量崩溃——在恐慌状态下（市场下跌后、波动率高企）伴随市场反弹，动量组合遭受罕见但持久的亏损，1927--2013 样本中最剧烈的两段是 1932 年 7--8 月与 2009 年 3--5 月\cite{danielmoskowitz2016}。注意：\emph{动量崩溃是趋势策略的“负偏”证据}——这提醒我们，趋势策略的偏度是“总体右偏、但尾部有深度左尾”，不是纯粹的右偏。

反转一侧的两段历史更直白。1998 年 LTCM 的收敛交易在俄罗斯违约后崩盘，其杠杆与流动性结构把一次“价差暂时扩大”放大为系统性失败\cite{lowenstein2000,jorion2000}；2007 年 8 月的“量化危机”中，因子型反转/统计套利组合被迫去杠杆，Khandani \& Lo(2011) 用因子模拟与 TAQ 交易数据识别出渐进的平仓过程，而非单日事件\cite{khandanilo2011}。两次事件的共同点是：反转策略的亏损不是渐进的，而是由\emph{不收敛+资金约束}共同触发的跳变。

\textbf{工程映射}。偏度差异决定了风险控制的手段不同。趋势策略的风险控制核心是\emph{波动率目标化}：Barroso \& Santa-Clara(2015) 用自身已实现方差缩放动量仓位，把夏普从 0.53 提升到 0.97、超额峰度从 18.2 降到 2.7\cite{barroso2015}；Harvey 等(2018) 在 60 多个资产上系统验证波动率目标改善了风险资产的夏普\cite{harvey2018}。反转策略的风险控制则不能依赖波动率缩放——理由见对偶四——它更依赖\emph{止损与资金冗余}：Leung \& Li(2015) 把止损退出纳入 OU 最优交易框架，正是因为不收敛的尾部风险无法靠仓位缩放消除\cite{leungli2015}。

\textbf{边界}。把趋势简单等同于“正偏”是危险的。动量崩溃的存在说明趋势策略也有深度左尾；而短期反转（提供流动性）的偏度未必是负的，因为流动性补偿本身可能呈正偏。偏度对偶严格成立于\emph{中低频、跨资产的组合口径}，在单资产高频繁口径上需要重新检验。

\subsection{对偶三：时间尺度——同一个自相关函数的三个区间}

\textbf{命题}：趋势与反转的归属由持有期决定，而非由信念决定。

\textbf{论证}。第~\ref{sec:unity}~节已经给出方差比 $VR(q)$ 的统一表述。这里补两点关键细节。

第一，美国市场的经验形状是\emph{波动}而非单调的。Lo \& MacKinlay(1988) 对 1962--1985 年 CRSP 数据做方差比检验，发现等权指数在周频显著偏离随机游走，而单只个股更接近随机游走\cite{lomackinlay1988}——这意味着反转信号在\emph{组合层面}比在\emph{个股层面}更容易被检验到，因为个股的特质噪声稀释了自相关。这是配对交易与截面反转选择“篮子”而非单资产的技术原因。

第二，长周期反转的证据本身存在争议。Richardson(1993) 明确质疑 Fama--French 的暂时成分结论，指出其证据对小样本与检验设定高度敏感\cite{richardson1993}。因此“长周期反转”的置信度明显低于“中周期动量”，这一点将在结论置信度章（第~\ref{sec:confidence}~节）体现。

\textbf{证据}。见第~\ref{sec:matrix}~节的统一比较矩阵中“持有期”一行，以及本文第~\ref{sec:empirical}~节的对照回测：同一资产池上，趋势与反转的相对夏普应当随 $L$ 单调反转。

\textbf{工程映射}。这一对偶有一个常被忽视的推论：\emph{“策略选择”其实可以转化为“持仓周期选择”}。组合管理的自由度因此从“选哪个策略”变成了“在哪几个尺度上分配风险预算”。Wood 等(2021) 的“慢动量加快反转”正是这一思路的实现：用同一信号的高频分量做反转、低频分量做趋势\cite{wood2021}。这不是把两个策略拼起来，而是承认它们是同一信号的两个频段。

\textbf{边界}。尺度依赖不是普适常数。交易成本、市场微观结构、投资者构成都会改变自相关函数的位置与幅度；同一 $L$ 在不同市场可能落在不同区间。因此“$L=12$ 个月是动量”这类说法是特定样本的经验概括，不是结构性定律。

\subsection{对偶四：风险来源——承担收敛风险 vs 收获延迟反应}

\textbf{命题}：反转策略的收益主要来自\emph{提供流动性、承担收敛失败风险}；趋势策略的收益主要来自\emph{风险溢价与信息延迟扩散}。前者是卖方（被资助方），后者是买方（资助方）。

\textbf{论证}。反转策略的经济角色是“当别人急着卖时买入、急着买时卖出”。这个动作本质上是做市：交易对手的紧迫需求把价格推离合理水平，反转交易者提供即时性，赚取补偿。其代价是承担两件事——价差可能继续扩大（不收敛），以及自己可能在需要资金时被抽资。Shleifer \& Vishny(1997) 的套利限制理论把第二点讲得最清楚：由于业绩与资金流的正反馈，套利者在错误定价加深时反而被迫减仓，从而无法等到收敛\cite{shleifervishny1997}。Brunnermeier \& Pedersen(2009) 的融资流动性模型给出同构的机制链条：价格下跌→保证金上升→被迫平仓→价格进一步下跌\cite{brunnermeier2009}。

趋势策略的经济角色是“在价格上涨后继续买入”。它的收益来源在文献中有两条解释。一条是风险补偿：趋势在极端市场中获利，说明它承担了某种在危机中才兑现的、被市场厌恶的风险\cite{hop2017}；另一条是行为：投资者对新信息的反应存在延迟，价格的调整不是一次完成，而是分步扩散。Hong \& Stein(1999) 的信息渐进扩散模型为后者提供了严格表述\cite{hongstein1999}，Barberis 等(1998) 与 Daniel 等(1998) 则从保守主义与过度自信两个角度给出互补解释\cite{bsv1998,dhs1998}。

\textbf{工程映射}。风险来源的差异解释了为什么两策略的“加仓纪律”方向相反。反转策略在亏损时应\emph{减少}而非增加敞口——因为亏损意味着不收敛的概率在上升，而资金约束会紧随其后。趋势策略在盈利时加仓是合理的（凸性），但在亏损时的加仓需要严格的波动率边界，否则会滑向动量崩溃。

\textbf{边界}。风险来源的二分不是互斥的。Ehsani \& Linnainmaa(2022) 的因子动量证据表明动量可以在因子层面被重新解释\cite{ehsani2022}；而反转收益中相当一部分被证明是流动性成本而非独立 alpha（Lo \& MacKinlay(1990) 的交叉自相关批评正是此意）\cite{lomackinlay1990}。因此更稳妥的表述是：\emph{两者的风险来源有主导倾向，而非排他归属}。

\subsection{对偶五：容量与换手——成本敏感度相差一个量级}

\textbf{命题}：反转策略的换手量级显著高于趋势策略，因而对交易成本敏感得多；趋势策略虽换手较低，却面临更严重的容量与拥挤问题。

\textbf{论证}。反转信号的半衰期短。以 OU 框架为例，最优交易带宽度与成本呈三次方根关系（Martin \& Sch\"oneborn 2011\cite{martin2011}），成本越高，带宽越宽、交易越稀——这意味着在给定成本下，反转策略被迫放弃一部分信号。Do \& Faff(2012) 的实证直接验证了这一点：把 Gatev 等的样本延伸并引入现实成本后，配对交易利润大幅衰减\cite{dofaff2012}。Korajczyk \& Sadka(2004) 用日内数据估计动量策略的成本，给出了盈亏平衡的基金规模\cite{korajczyksadka2004}；Frazzini 等(2018) 用实盘交易数据估计异象的真实成本，指出换手高的因子净收益衰减最重\cite{frazzini2018}；Novy-Marx \& Velikov(2016) 建立的异象成本分类法也指向同一结论\cite{novymarxvelikov2016}。

趋势一侧，King 等(2016)\cite{kim2016}与 Baltas \& Kosowski(2013) 的研究显示，时间序列动量的换手可以通过信号平滑（如线性趋势拟合、Yang--Zhang 区间估计量）显著降低，且样本外表现更好。换手低意味着成本敏感度低——这是趋势策略在长期样本中“看似稳健”的一个技术原因。

但趋势的代价是拥挤。Stein(2009) 的理论指出，对于缺乏锚定的策略（动量是典型代表），套利者无法知道有多少同行在同时进场，协调失败会放大价格偏离；同时杠杆外溢制造火售外部性\cite{stein2009}。Lou \& Polk(2022) 提出的 comomentum 指标从相关性角度度量动量拥挤度，发现高拥挤之后动量收益下降\cite{loupolk2022}；Baltas(2019) 系统研究了另类风险溢价策略的拥挤与容量\cite{baltas2019}。

\textbf{工程映射}。这一对偶决定了两个策略的规模天花板不同。反转策略的约束来自\emph{成本}：单笔交易的期望收益在扣除成本后可能所剩无几，因此它天然适合规模较小、速度较快的资金。趋势策略的约束来自\emph{拥挤}：信号本身在低换手下可容纳大资金，但同质化导致信号失效，因此它天然适合规模较大、但需要持续分散与容量监控的资金。

\textbf{边界}。这个二分有例外：跨资产的趋势策略在流动性差的品种上换手成本同样很高；而某些低频反转（如多年尺度的价值反转）换手极低。因此准确的表述应绑定\emph{尺度}：短尺度反转成本敏感，中尺度趋势拥挤敏感。

\subsection{对偶六：失效模式——“慢慢失血”与“突然跳崖”的错配}

\textbf{命题}：两个策略都会亏钱，但亏法不同。反转策略的失效是“不收敛+资金断裂”的联合事件，表现为突然的深亏；趋势策略的失效是“震荡市慢性出血”叠加“崩溃时跳崖”。

\textbf{论证}。反转策略的典型失败路径在两次历史事件中被完整记录。LTCM 的收敛交易在价差扩大时本应加仓等待，但杠杆与赎回约束让它无法等到收敛\cite{lowenstein2000,jorion2000}。2007 年 8 月的量化危机中，多个统计套利基金同时去杠杆，Khandani \& Lo(2011) 识别出的两次平仓事件集中在数小时到数日内完成，起始头寸集中于金融股、多账面市值比、空盈利动量\cite{khandanilo2011}。共同的结构是：\emph{亏损本身不可怕，可怕的是亏损触发了资金约束}。

趋势策略的失败有两副面孔。震荡市中是慢性出血：价格反复穿越信号阈值，策略不断止损，收益为负但幅度温和，波动率目标化会进一步降低敞口（Harvey 等 2018 指出这正是波动率目标的代价之一\cite{harvey2018}）。危机中是突然跳崖：动量崩溃。Daniel \& Moskowitz(2016) 强调崩溃的\emph{可预测性}——它发生在恐慌状态，且部分可由市场波动率与市场收益预测，据此构造的动态动量策略把 alpha 与夏普大约翻倍\cite{danielmoskowitz2016}。

\textbf{工程映射}。失效模式差异对应两套不同的“熔断”设计。反转策略需要的是\emph{资金冗余与止损线}：设定价差偏离的容忍上限，超过则减仓，且预留足够的保证金缓冲以防被迫平仓；不能用“越跌越买”的加仓逻辑。趋势策略需要的是\emph{状态过滤}：在恐慌状态下降低甚至反转敞口，避免在崩溃中继续追涨。

\textbf{边界}。两策略的失效在极端去杠杆事件中会同时发生——2007 年 8 月中，量化反转基金与动量基金都被冲击，因为冲击来自共同的资金端而非各自的信号端\cite{khandanilo2011,brunnermeier2009}。这说明\emph{分散化在资金冲击下会失效}，这是组合层面而非策略层面的问题。

\subsection{对偶七：入场门槛——先验证后交易 vs 边交易边估计}

\textbf{命题}：反转策略在结构上多一道“资格检验”；趋势策略没有这道门槛，但需要持续估计信号的稳定性。前者是“条件性”策略，后者是“无偏条件”策略。

\textbf{论证}。配对交易的标准流程是：选对（距离法或协整法）→ 检验价差平稳性 → 估计 OU 参数（回归速度、半衰期）→ 设定阈值带 → 交易。每一步都是筛选，不合格的样本被排除。Vidyamurthy(2004) 的专著系统化了这一流程\cite{vidyamurthy2004}；Elliott 等(2005) 用 Kalman 滤波在线估计 OU 状态，把“资格检验”从一次性改为滚动进行\cite{elliott2005}。

趋势策略的流程是：估计动量统计量（过去收益的符号、排序或回归斜率）→ 按波动率缩放建仓 → 定期调仓。没有独立的资格检验步骤，因为动量信号对每一个资产都可以计算，不需要先证明该资产“适合动量”。

这两条流程的差别带来可检验的推论：\emph{反转策略的有效样本数少于趋势策略}。反转需要先筛选出满足条件的配对/篮子，筛选失败的资产被浪费；趋势则可以用满整个资产池。在资产数量有限的市场（如单一期货品种），反转策略的样本基础更薄。

\textbf{工程映射}。这道门槛使反转策略的\emph{超参数}更多（选对规则、协整显著性、半衰期上下限、阈值带宽度），因而数据窥探的空间更大。这既是它容易过拟合的原因，也解释了为何配对交易文献中“利润衰减”证据如此普遍——Do \& Faff(2012) 的发现可以部分归因于早期研究的参数选择在样本外失效\cite{dofaff2012}。

\textbf{边界}。这个二分有反例：截面动量的“资格”隐含在排序中（只有流动性足够的资产才能进入组合），所以趋势策略并非完全没有筛选，只是筛选维度不同（流动性 vs 平稳性）。

\subsection{旁支落地：配对筛选的七个步骤}

对偶七断言反转多一道资格检验。把这道检验写成算法，就是配对交易文献的标准流程\cite{vidyamurthy2004}——七步，每步一个判据，不合格的样本在原地被拒绝，走完全程的才进入交易。

\begin{enumerate}[leftmargin=2.2em]
  \item \textbf{分组生成候选对}。只在同产业链或同板块内两两组合。跨组配对缺少经济意义上的长期均衡，协整检验在这种对上很容易给出伪通过。
  \item \textbf{相关性初筛}。形成期两序列对数价格的相关系数设下限（常见取 $0.6$）。这一步不是统计检验，是算力与稳健性的双重过滤——低相关配对即便偶然通过后续检验，价差波动也大到不可交易。
  \item \textbf{协整检验}。Engle--Granger 两步法\cite{englegranger1987}：对形成期做回归 $\ln P^A_t=\alpha+\beta\ln P^B_t+e_t$，再对残差 $e_t$ 做单位根检验\cite{dickeyfuller1979}，$p<0.05$ 才放行。这是全流程唯一的“资格检验”，对应对偶七命题里的平稳性要求。
  \item \textbf{回复速度估计}。对残差拟合稳定 AR(1)：$e_{t+1}=c+\phi e_t+u_t$，要求 $0<\phi<1$；换算连续时间参数 $\theta=-\ln\phi/\Delta$，半衰期 $H=\ln 2/\theta$。$\phi\ge1$ 意味着价差不回归，$\phi\le0$ 意味着高频振荡而非缓慢回复，两者都让半衰期失去意义。
  \item \textbf{半衰期区间筛选}。$H$ 设上下界：下界挡掉回复过快的对（交易频率高、成本吃光价差），上界挡掉占用资金过久的对（期间结构性断裂的概率随时间上升）。
  \item \textbf{滚动重估}。形成期与交易期严格分离、滚动推进，杜绝用交易期数据选对造成的 look-ahead——这一范式来自 Gatev 等(2006) 的 $6+6$ 个月设计\cite{gatev2006}；交易期的均值与标准差须用滚动统计量，不能沿用形成期常数，否则波动放大时 $|z|$ 虚高、开仓即触止损。
  \item \textbf{生成可交易对表}。输出对冲比率（静态回归系数，或 Kalman 滤波的在线估计\cite{elliott2005}）、价差均值与标准差、半衰期，据此设定开平仓阈值带。
\end{enumerate}

七步的顺序不可调换，且每一步的参数（相关性下限、显著性水平、半衰期区间、阈值带宽）都是自由度。这一点必须说透：\emph{筛选算法本身也在做多重检验}——一百个候选对过筛后剩五个“平稳”对，这五个的平稳性是选择的结果，不是证据。因此筛选参数必须在实验前锁定、并整体计入第~\ref{sec:audit}~节的试验族审计，否则筛选就从资格检验退化为变相的数据窥探。

\subsection{七组对偶的合并判断}

把七组对偶合起来，得到一个修正后的判断：

\begin{quote}
\itshape
趋势与均值回归共享同一个信号，却生活在两套约束里。
反转策略在一个“需要资格证明”的世界里做流动性生意，它的敌人是成本与资金约束，
它的失效方式是突然的；趋势策略在一个“无需资格证明”的世界里做延迟反应的生意，
它的敌人是拥挤与状态误判，它的失效方式是两面性的——慢性出血与尾部跳崖。
\end{quote}

这个判断是本文的第二个核心结论。它比“两策略对立”更精确，也比“两策略等价”更诚实。第~\ref{sec:execution}~节将把它翻译成具体的执行差异清单。
